\chapter{The KV Cache}
\label{ch:kvcache}

The previous chapter showed that every token generated by a transformer requires attending to all previous tokens---computing query-key dot products and value-weighted sums across the entire context. Without intervention, this means re-running the full attention computation from scratch for every single new token. The \emph{KV cache} is the data structure that eliminates this redundancy: it stores the key and value vectors from previous positions so they never need to be recomputed. This chapter explains why the KV cache exists, what it looks like in memory, why it dominates GPU resource consumption, and how two critical innovations---PagedAttention and prefix sharing---make it manageable.

\section{Why Cache?}
\label{sec:kvcache:why}

\marginnote{\emph{Autoregressive generation:} producing one token at a time, each conditioned on all previous tokens.}%
Recall from Chapter~\ref{ch:transformer} that autoregressive generation produces tokens one at a time. At step $t$, the model has already generated positions $1, \ldots, t$ and must compute attention over all of them to produce token $t+1$. The attention mechanism requires key and value vectors for every position in the sequence. The question is: where do those vectors come from?

\subsection{The Na\"ive Approach}

To understand the problem, we need to distinguish two phases of inference. \textbf{Prefill} is the initial phase: the user's entire prompt arrives at once, and the model processes all $n$ prompt tokens in a single parallel forward pass. This is efficient---large batches of matrix multiplications keep the GPU busy. \textbf{Decode} is what comes next: the model generates tokens one at a time, each conditioned on every token before it. Decode is where the waste lives.

Here is why. When the model generates token $t+1$ during decode, it needs to compute attention over positions $1$ through $t$. Attention requires a \emph{key} and \emph{value} vector for every one of those positions (these are the $\mathbf{K}$ and $\mathbf{V}$ matrices from Chapter~\ref{ch:transformer}). Without a cache, the model has no memory of previous steps---so it must recompute every key and value vector from scratch, by re-running the projection $\mathbf{k}_j = \mathbf{W}_K \mathbf{x}_j$ and $\mathbf{v}_j = \mathbf{W}_V \mathbf{x}_j$ for every past position $j$, at every layer. Then for token $t+2$, it does the same thing for positions $1$ through $t+1$---and all the key/value work from the previous step is thrown away and redone.

The waste compounds: position 1's key and value vectors are recomputed at every single decode step. Position 2's vectors are recomputed at every step after the second. And so on. The model is doing the same matrix multiplications over and over, producing the same results each time, because it has nowhere to store them between steps.

Let us quantify the waste. Suppose we start with a prompt of $n$ tokens and want to generate $T$ new tokens. At decode step $i$ (for $i = 1, \ldots, T$), the model must compute key and value projections for all $n + i - 1$ previous positions.\footnote{The projection for position $j$ is $\mathbf{k}_j = \mathbf{W}_K \mathbf{x}_j$ and $\mathbf{v}_j = \mathbf{W}_V \mathbf{x}_j$, each costing $O(d^2)$ operations per layer per head. We absorb the per-element cost into a constant and focus on the scaling in $n$, $T$, and $L$.} The total number of key-value projections across all decode steps is:
\begin{equation}
\label{eq:naive-projections}
\sum_{i=1}^{T} (n + i - 1) = nT + \frac{T(T-1)}{2} = O(nT + T^2)
\end{equation}
Each projection is computed at every layer, so the total work scales as $O\bigl(L \cdot (nT + T^2) \cdot d^2\bigr)$, where $L$ is the number of layers and $d$ is the model dimension. For long generations where $T$ is large, the $T^2$ term dominates---the system does \emph{quadratic} work in the number of generated tokens, not because attention itself is quadratic (that is a separate cost), but because it redundantly recomputes projections.

\subsection{The Caching Insight}

\marginnote{\emph{Immutability:} the key and value vectors for position $j$ depend only on the input at position $j$ and the frozen model weights. They are the same at every decode step.}%
The key observation is simple: the key and value vectors for positions $1, \ldots, t-1$ do not change when we generate token $t$. They depend only on the input at their own position and the (frozen) model weights. We can compute them once, store them, and never recompute them.

With a cache, each decode step computes only the new token's key and value projections---a single row per layer per head---appends them to the stored cache, and performs attention against the full cached sequence. The total number of key-value projections drops from Equation~\ref{eq:naive-projections} to:
\begin{equation}
\label{eq:cached-projections}
n + T
\end{equation}
where $n$ accounts for the initial prompt processing (prefill) and $T$ for the $T$ decode steps, each contributing exactly one new token's projections. Per layer, the projection cost becomes $O\bigl(L \cdot (n + T) \cdot d^2\bigr)$---\emph{linear} in the total sequence length.

\begin{insightbox}
The KV cache trades memory for compute. Without it, each decode step recomputes all keys and values---quadratic total work over a generation. With it, each step computes only the new token's key and value vectors, appends them to the cache, and performs attention against the cached entries. The total projection work drops from $O(nT + T^2)$ to $O(n + T)$ --- a factor-of-$T$ reduction for the projection cost. The attention computation itself ($O(t \cdot d)$ per step, $O(T^2 d / 2)$ total) is inherent and irreducible without approximation; the cache eliminates only the \emph{redundant} projection work, which is the part that was entirely wasted.
\end{insightbox}

\subsection{A Concrete Example}

Consider serving LLaMA-70B with a prompt of $n = 2{,}048$ tokens and a generation length of $T = 512$. Without the cache, Equation~\ref{eq:naive-projections} gives:
\[
2{,}048 \times 512 + \frac{512 \times 511}{2} = 1{,}048{,}576 + 130{,}816 = 1{,}179{,}392 \text{ projection sets}
\]
With the cache, Equation~\ref{eq:cached-projections} gives $2{,}048 + 512 = 2{,}560$ projection sets. The cache reduces the total projection work by a factor of approximately $460 \times$. Each ``projection set'' involves computing a key and value vector across all 80 layers, so the absolute compute savings are enormous: millions of matrix-vector multiplications avoided.

The cost of this saving is memory. Those $2{,}560$ projection sets must be stored somewhere, and as we will see in Section~\ref{sec:kvcache:memory}, for large models this storage dominates GPU memory consumption.

\section{The KV Cache Data Structure}
\label{sec:kvcache:structure}

For each layer $\ell \in \{1, \ldots, L\}$ and each key-value head $h \in \{1, \ldots, H_{\text{KV}}\}$, the KV cache stores two growing tensors:
\begin{align}
\mathbf{K}^{(\ell,h)}_{\text{cache}} &\in \mathbb{R}^{t \times d_k} \\
\mathbf{V}^{(\ell,h)}_{\text{cache}} &\in \mathbb{R}^{t \times d_v}
\end{align}
where $t$ is the current sequence length (which grows by one at each decode step), $d_k$ is the key head dimension, and $d_v$ is the value head dimension (usually $d_k = d_v = d / H$ for the original multi-head attention, but we keep them separate for generality). At each decode step, the new token's key and value vectors are computed and appended as a new row to each cache tensor.

\subsection{Memory Layout}

The full KV cache for a single sequence can be viewed as a 5-dimensional tensor:
\[
\text{Cache} \in \mathbb{R}^{2 \times L \times H_{\text{KV}} \times t \times d_k}
\]
where the leading dimension of 2 accounts for the key and value tensors. In practice, the physical layout varies by implementation:

\begin{itemize}
  \item \textbf{Layer-first layout:} $[L, 2, H_{\text{KV}}, t, d_k]$. Each layer's cache is contiguous, which is natural because attention is computed one layer at a time during the forward pass. This is the layout used by Hugging Face Transformers.
  \item \textbf{Sequence-first layout:} $[2, t, L, H_{\text{KV}}, d_k]$. Each token's full cache across all layers is contiguous, which can improve cache locality for certain memory management strategies.
  \item \textbf{Block layout:} $[2, L, H_{\text{KV}}, \lceil t / B \rceil, B, d_k]$. The sequence dimension is divided into fixed-size blocks of $B$ tokens. This is the layout used by PagedAttention (Section~\ref{sec:kvcache:paged}).
\end{itemize}

\marginnote{\emph{GQA (Grouped-Query Attention):} multiple query heads share a single key-value head, reducing cache size by the group factor $G = H / H_{\text{KV}}$.}%
The crucial parameter is $H_{\text{KV}}$, the number of key-value heads. In the original transformer, every query head has its own key and value head: $H_{\text{KV}} = H$. Modern models use \emph{grouped-query attention} (GQA),\footnote{Joshua Ainslie, James Lee-Thorp, Michiel de Jong, Yinfei Yang, Siddhartha Reddy Jonnalagadda, and Santiago Onta\~{n}\'{o}n, ``GQA: Training Generalized Multi-Query Transformers from Multi-Head Checkpoints,'' \emph{EMNLP}, 2023.} where a group of $G$ query heads shares a single KV head, so $H_{\text{KV}} = H / G$. At the extreme, \emph{multi-query attention} (MQA) uses $H_{\text{KV}} = 1$: all query heads share a single key-value head.

GQA is a direct optimization of cache size. If a model has $H = 64$ query heads but only $H_{\text{KV}} = 8$ KV heads (a group factor of $G = 8$), the KV cache is $8\times$ smaller than it would be with standard multi-head attention. This is not an approximation trick; models are explicitly \emph{trained} with GQA, learning to share KV representations across groups of query heads. LLaMA~2~70B uses $H = 64$ query heads and $H_{\text{KV}} = 8$ GQA heads; LLaMA~3 continues this pattern.

\subsection{Growth Dynamics}

\marginnote{\emph{Prefill:} the initial phase where the entire prompt is processed in one parallel forward pass.\\[4pt]\emph{Decode:} the autoregressive phase where tokens are generated one at a time.}%
During \emph{prefill}, the cache is populated in one shot: the entire prompt of $n$ tokens is processed in parallel, and all $n$ key-value pairs are computed simultaneously at every layer. The cache goes from empty to $t = n$ in a single forward pass.

During \emph{decode}, the cache grows by exactly one row per step. At step $i$, the cache size is $t = n + i$, and the new token's key/value pair is written to position $t$ in each layer's cache. The attention computation reads all $t$ rows to produce the output for the new token.

This asymmetry---bulk write during prefill, incremental append during decode---has profound implications for system design:

\begin{itemize}
  \item \textbf{Prefill is compute-bound.} Processing $n$ tokens in parallel involves large matrix multiplications ($n \times d$ matrices), which saturate the GPU's compute units. The arithmetic intensity (FLOPs per byte of memory accessed) is high.
  \item \textbf{Decode is memory-bound.} Processing 1 token involves reading the entire KV cache ($t$ rows per layer per head) to compute a single dot product per row. The arithmetic intensity is low---most of the time is spent waiting for data to arrive from HBM, not computing.
\end{itemize}

This fundamental asymmetry means that prefill and decode have \emph{opposing} hardware requirements: prefill wants more compute (more FLOPS), while decode wants more memory bandwidth (more HBM bandwidth and capacity). Running both on the same GPU forces a compromise. This is why Dynamo supports \emph{disaggregated serving}: running prefill and decode on separate, independently scaled GPU pools (Chapter~\ref{ch:distributed}). But disaggregation introduces a new problem: after prefill computes the KV cache on a prefill GPU, it must be \emph{transferred} to a decode GPU. The KV cache, which we have just seen can be gigabytes in size, must cross GPU-to-GPU interconnects---making efficient cache transfer and placement a first-class system design concern.

\section{Memory Analysis}
\label{sec:kvcache:memory}

We have seen that the KV cache eliminates redundant computation at the cost of storing every past token's keys and values. How much memory does this actually require? The answer is sobering: the KV cache is the dominant consumer of GPU memory during inference, often exceeding the model weights themselves for large batches with long sequences. Let us build up the calculation from first principles.

\subsection{Per-Token Cache Size}

For a single token at a single layer and a single KV head, the cache stores one key vector ($d_k$ elements) and one value vector ($d_v$ elements). With $d_k = d_v = d_{\text{head}}$ and using $b$ bytes per element (e.g., $b = 2$ for FP16, $b = 1$ for INT8):
\begin{equation}
\label{eq:per-token-per-layer-per-head}
\text{Bytes per token per layer per head} = 2 \times d_{\text{head}} \times b
\end{equation}
Summing across all $H_{\text{KV}}$ heads and $L$ layers:
\begin{equation}
\label{eq:per-token}
\text{Bytes per token} = 2 \times L \times H_{\text{KV}} \times d_{\text{head}} \times b
\end{equation}

\begin{notebox}
\textbf{KV cache size for LLaMA-70B.} With $L = 80$ layers, $H_{\text{KV}} = 8$ GQA heads, $d_{\text{head}} = 128$, and FP16 ($b = 2$ bytes per element):
\[
\text{Cache per token} = 2 \times 80 \times 8 \times 128 \times 2 = 327{,}680\;\text{bytes} \approx 320\;\text{KB}
\]
Every single token in every active sequence consumes 320\,KB of GPU memory just for its cached keys and values. For context, this is larger than the entire L1 cache on most CPUs.
\end{notebox}

Contrast this with a smaller model. For LLaMA-7B ($L = 32$, $H_{\text{KV}} = 32$ with standard MHA, $d_{\text{head}} = 128$, FP16):
\[
\text{Cache per token} = 2 \times 32 \times 32 \times 128 \times 2 = 524{,}288\;\text{bytes} \approx 512\;\text{KB}
\]
The 7B model actually has a \emph{larger} per-token cache than the 70B model because it uses standard multi-head attention ($H_{\text{KV}} = 32$) rather than GQA ($H_{\text{KV}} = 8$). This illustrates why GQA was adopted: it reduces cache size by the group factor without proportional loss in model quality.

\subsection{Batch-Level Memory}

\marginnote{\emph{Batch size:} the number of sequences being processed concurrently. Larger batches amortize the cost of loading model weights but require proportionally more KV cache memory.}%
For a batch of $B$ sequences, each of length $s$ (including both prompt and generated tokens so far):
\begin{equation}
\label{eq:batch-cache}
\text{Total cache memory} = B \times s \times 2 \times L \times H_{\text{KV}} \times d_{\text{head}} \times b
\end{equation}

\begin{notebox}
\textbf{Concrete scenario: LLaMA-70B on an A100-80GB.} The model weights in FP16 occupy approximately 140\,GB, requiring at least 2 A100-80GB GPUs with tensor parallelism. Suppose we split across 2 GPUs, so each GPU holds about 70\,GB of weights. That leaves roughly 10\,GB per GPU for the KV cache.\footnote{In practice, the CUDA runtime, activation memory, and other overheads consume additional memory, leaving less than 10\,GB.}

With 320\,KB per token, 10\,GB supports:
\[
\frac{10 \times 1024^3}{320 \times 1024} \approx 32{,}768 \text{ token-slots}
\]
If each sequence uses up to 4,096 tokens, we can serve at most $32{,}768 / 4{,}096 = 8$ concurrent sequences per GPU. With 2,048-token sequences, we get 16. These are small batches---and batch size directly determines throughput, because larger batches amortize the cost of loading model weights from HBM.
\end{notebox}

This reveals the fundamental tension: \textbf{throughput requires large batches, but large batches require large KV caches, and GPU memory is fixed.} A system that cannot manage cache memory efficiently will either reject requests (reducing throughput) or run out of memory and crash. This is the problem that PagedAttention solves.

\begin{whynotbox}{Why not just add more GPU memory?}
A natural reaction to the KV cache memory problem is: buy bigger GPUs. If 80\,GB is not enough, wait for 160\,GB or 320\,GB parts. This does not solve the problem---it merely shifts the bottleneck.

\textbf{Memory grows linearly; demand grows combinatorially.} Doubling HBM capacity lets you double the batch size \emph{or} double the sequence length---but real workloads want both. A model serving 128K-token contexts at batch size 64 needs $128{,}000 \times 64 \times 320\;\text{KB} \approx 2.5$\,TB of cache for LLaMA-70B. No single GPU comes close.

\textbf{Cost.} HBM is the most expensive component on a GPU die. The A100-80GB costs roughly 2$\times$ the A100-40GB, and the H100-141GB is priced accordingly. Doubling memory to double batch size is a poor trade-off when software can achieve a 2--5$\times$ improvement through better allocation (PagedAttention) at zero hardware cost.

\textbf{Bandwidth, not just capacity.} Even if a GPU had unlimited memory, the decode phase is bottlenecked by \emph{memory bandwidth}---how fast data can be read from HBM, not how much HBM exists. Adding capacity without proportionally increasing bandwidth yields diminishing returns.

The right approach is to use the memory you have more efficiently. That is precisely what PagedAttention and prefix sharing accomplish.
\end{whynotbox}

\subsection{KV Cache vs.\ Model Weights}

For a single short request, the KV cache is much smaller than the model weights. But at batch scale, the cache grows while the weights stay constant:
\begin{center}
\begin{tabular}{lrr}
\hline
\textbf{Component} & \textbf{LLaMA-7B (FP16)} & \textbf{LLaMA-70B (FP16)} \\
\hline
Model weights & 14 GB & 140 GB \\
KV cache (1 seq, 2K tokens) & 1.0 GB & 0.625 GB \\
KV cache (1 seq, 4K tokens) & 2.0 GB & 1.25 GB \\
KV cache (32 seq, 4K tokens) & 64 GB & 40 GB \\
KV cache (64 seq, 4K tokens) & 128 GB & 80 GB \\
\hline
\end{tabular}
\end{center}

At a batch size of 32 with 4K-token sequences, the KV cache for LLaMA-7B ($\sim$64\,GB) is over $4\times$ the model weights (14\,GB). For LLaMA-70B, the 40\,GB cache is a smaller fraction of the 140\,GB weights, but still represents an enormous memory demand that must be carefully managed.

Two trends in model design work in opposite directions. \textbf{GQA reduces cache pressure}: by sharing KV heads across groups of query heads, GQA cuts the per-token cache cost proportionally. LLaMA-70B's use of 8 KV heads instead of 64 reduces its per-token cache from $\sim$2.6\,MB (if it used standard MHA) to $\sim$320\,KB---an $8\times$ reduction. \textbf{Longer context windows increase cache pressure}: extending the maximum context from 4K to 128K tokens multiplies the worst-case cache size by $32\times$. Models like LLaMA~3 support 128K-token contexts, meaning a single sequence could consume $128{,}000 \times 320\;\text{KB} \approx 40$\,GB of cache---more than the entire HBM of many GPUs. These opposing forces make KV cache management one of the most active areas of systems research in LLM inference.

\begin{insightbox}
The KV cache is the reason batch size is limited by GPU memory, not by model weights. Model weights are a fixed cost; the KV cache scales linearly with both the number of concurrent sequences and their lengths. Any system that wants to maximize throughput must minimize KV cache waste---and the dominant source of waste, before PagedAttention, was memory fragmentation from na\"ive allocation strategies.
\end{insightbox}

\section{PagedAttention}
\label{sec:kvcache:paged}

\marginnote{\emph{PagedAttention:} applies OS-style virtual memory paging to KV cache management, eliminating fragmentation. Introduced in the vLLM paper by Kwon et al., 2023.}%
Having established that the KV cache dominates GPU memory and that its size grows linearly with batch size and sequence length, we now turn to the question of \emph{how} that memory is managed. The na\"{\i}ve approach wastes most of it. PagedAttention reclaims that waste.

Traditional KV cache implementations pre-allocate a contiguous block of memory for each sequence at its maximum possible length. If the system's maximum sequence length is 4,096 tokens, every incoming request reserves $4{,}096 \times 320$\,KB $= 1.28$\,GB of cache memory (for LLaMA-70B)---even if the actual generation produces only 50 tokens. The waste is severe: Kwon et al.\footnote{Woosuk Kwon, Zhuohan Li, Siyuan Zhuang, Ying Sheng, Lianmin Zheng, Cody Hao Yu, Joseph E.\ Gonzalez, Hao Zhang, and Ion Stoica, ``Efficient Memory Management for Large Language Model Serving with PagedAttention,'' \emph{SOSP}, 2023.} measured that existing systems waste 60--80\% of KV cache memory due to three forms of fragmentation:

\marginnote{\emph{Fragmentation:} wasted memory due to allocation policies---space that is allocated but unused (internal) or free but unusable (external).}%
\begin{enumerate}
  \item \textbf{Reservation waste.} Space is reserved for the maximum sequence length but most sequences are shorter. A 50-token response in a 4,096-token slot wastes 98.8\% of its allocation.
  \item \textbf{Internal fragmentation.} Even for completed sequences, the allocated buffer may be larger than the final sequence length due to alignment or allocation granularity.
  \item \textbf{External fragmentation.} As sequences arrive and depart at different times, the free memory becomes scattered into non-contiguous chunks, each too small to serve a new request even though the total free memory is sufficient.
\end{enumerate}

\subsection{The Paging Analogy}

PagedAttention applies the same insight that operating systems use for virtual memory: \textbf{decouple the logical layout from the physical layout}. Instead of allocating one contiguous block per sequence, PagedAttention divides GPU memory into fixed-size \emph{pages} (called \emph{blocks} in vLLM's terminology),\footnote{We use ``block'' and ``page'' interchangeably in this chapter. vLLM's codebase uses ``block''; the conceptual analogy is to OS pages.} each holding the keys and values for a fixed number of tokens.

\begin{notebox}
\textbf{Block size.} A typical block size is $B = 16$ tokens. For LLaMA-70B in FP16, one block stores:
\[
16 \times 320\;\text{KB} = 5{,}120\;\text{KB} = 5\;\text{MB}
\]
per layer-head slice. The total block size across all layers and heads is 5\,MB. A single A100-80GB can hold roughly $80{,}000 / 5 \approx 16{,}000$ blocks (in the memory left after model weights), tracking them via a simple free list.
\end{notebox}

\begin{whynotbox}{Why fixed-size blocks instead of variable-size allocations?}
Operating systems faced this same design choice decades ago. Early memory allocators used variable-size allocations, which leads to \emph{external fragmentation}: after many allocations and frees, memory is littered with differently-sized holes, and a request for $N$ bytes can fail even when the total free memory exceeds $N$.

Fixed-size blocks eliminate external fragmentation entirely. Any freed block can satisfy any allocation request---there is no need to search for a ``big enough'' hole or compact memory. The trade-off is \emph{internal fragmentation}: the last block of a sequence is typically partially filled. With block size $B = 16$, the expected waste is $B/2 = 8$ tokens per sequence, which for LLaMA-70B amounts to $8 \times 320\;\text{KB} = 2.5$\,MB---negligible compared to the tens of gigabytes that reservation-based allocation wastes.

A variable-size scheme could waste zero bytes at the tail, but would reintroduce the fragmentation nightmare that makes batch scheduling unpredictable. Fixed blocks give the allocator $O(1)$ allocation and deallocation (just push/pop from a free list), which matters when the scheduler must make allocation decisions on the critical path of every decode step.
\end{whynotbox}

Each sequence maintains a \emph{block table}: an array of pointers mapping logical block indices to physical block addresses. Position $j$ in the sequence maps to logical block $\lfloor j / B \rfloor$, offset $j \bmod B$ within that block. As the sequence grows during decode, new physical blocks are allocated from a shared pool on demand---only when the current block is full and a new token needs space.

\subsection{How It Works}

Consider a sequence that has generated 35 tokens with a block size of 16:
\begin{itemize}
  \item Logical block 0 maps to physical block 7 (tokens 0--15, full)
  \item Logical block 1 maps to physical block 42 (tokens 16--31, full)
  \item Logical block 2 maps to physical block 3 (tokens 32--34, partially filled)
\end{itemize}
Physical blocks 7, 42, and 3 are not contiguous in GPU memory---they could be anywhere. The block table provides the indirection. When token 36 arrives, it is appended to physical block 3 at offset 3. When physical block 3 becomes full (at token 47), a new physical block is allocated from the free pool for logical block 3.

\begin{figure}[ht]
\centering
\begin{tikzpicture}[
    blocktable/.style={rectangle, draw, minimum width=1.6cm, minimum height=0.6cm, align=center, font=\small},
    physblock/.style={rectangle, draw, minimum width=2.4cm, minimum height=0.6cm, align=center, font=\small},
    fullblock/.style={physblock, fill=blue!15},
    partialblock/.style={physblock, fill=blue!8, pattern=north east lines, pattern color=blue!15},
    freeblock/.style={physblock, fill=gray!10, dashed},
    arrow/.style={-{Stealth[length=5pt]}, thick},
    label/.style={font=\footnotesize\itshape},
]
    % Block table (left side)
    \node[font=\small\bfseries] at (0, 2.2) {Block Table (Seq.\ A)};
    \node[blocktable] (bt0) at (0, 1.5) {Logical 0};
    \node[blocktable] (bt1) at (0, 0.8) {Logical 1};
    \node[blocktable] (bt2) at (0, 0.1) {Logical 2};

    % Physical memory (right side) -- scattered
    \node[font=\small\bfseries] at (5.5, 3.0) {GPU Physical Memory};

    \node[freeblock] (p0) at (5.5, 2.3) {\textcolor{gray}{Free}};
    \node[freeblock] (p1) at (5.5, 1.6) {\textcolor{gray}{Free}};
    \node[fullblock]  (p3) at (5.5, 0.9) {Block 3: tok 32--34};
    \node[freeblock] (p4) at (5.5, 0.2) {\textcolor{gray}{Free}};
    \node[fullblock]  (p7) at (5.5, -0.5) {Block 7: tok 0--15};
    \node[freeblock] (p8) at (5.5, -1.2) {\textcolor{gray}{Free}};
    \node[fullblock] (p42) at (5.5, -1.9) {Block 42: tok 16--31};

    % Arrows from block table to physical blocks
    \draw[arrow, blue!70!black] (bt0.east) -- (p7.west) node[label, above, midway, sloped] {ptr = 7};
    \draw[arrow, blue!70!black] (bt1.east) -- (p42.west) node[label, below, midway, sloped] {ptr = 42};
    \draw[arrow, blue!70!black] (bt2.east) -- (p3.west) node[label, above, midway, sloped] {ptr = 3};

    % Annotation
    \node[label, text width=3.2cm, align=center] at (5.5, -2.8) {Physical blocks are scattered\\across GPU memory};
\end{tikzpicture}
\caption{PagedAttention block table for a sequence of 35 tokens with block size 16. The block table maps logical block indices to physical block addresses. Physical blocks need not be contiguous---the indirection layer decouples logical sequence order from physical memory layout, exactly as a page table does in an operating system.}
\label{fig:block-table}
\end{figure}

When the sequence finishes (e.g., the model generates an end-of-sequence token), all its physical blocks are returned to the free pool. Because blocks are fixed-size, any freed block can immediately serve any new sequence---there is no external fragmentation. And because blocks are allocated on demand rather than reserved up front, there is no reservation waste.

\subsection{The Block Table}

\marginnote{\emph{Block table:} a per-sequence array mapping logical block indices to physical block addresses in GPU memory---analogous to a page table in an OS.}%
The block table is the central bookkeeping structure. For each active sequence, the system maintains an array:
\[
\text{block\_table}[i] = \text{physical block index for logical block } i
\]
To access the KV cache for token at position $j$:
\begin{enumerate}
  \item Compute the logical block index: $\lfloor j / B \rfloor$
  \item Look up the physical block: $\text{block\_table}[\lfloor j / B \rfloor]$
  \item Compute the offset within the block: $j \bmod B$
  \item Read or write at that offset in the physical block
\end{enumerate}

The block table itself is small---a sequence of 4,096 tokens with block size 16 needs only $4{,}096 / 16 = 256$ entries, each a 32-bit pointer. The total block table overhead across hundreds of concurrent sequences is negligible compared to the cache data itself.

\marginnote{\emph{Reference counting:} tracking how many sequences point to each physical block. A block is freed only when its count reaches zero.}%
When a sequence finishes, its entire block table is scanned and each physical block's reference count is decremented. Blocks with a reference count of zero are returned to the free pool. This deallocation is $O(s / B)$ per sequence---typically a few hundred operations at most.

\subsection{Modified Attention Kernel}

The attention kernel must be modified to follow block table pointers instead of assuming contiguous memory. For each query position, the kernel iterates over the sequence's block table, fetches the key and value data from the indicated physical block, and performs the dot product and weighted sum. The mathematical computation is identical to standard attention---the outputs are bit-for-bit the same. PagedAttention is purely a memory management optimization, not an approximation.

The kernel modification introduces one additional memory indirection per block (reading the block table entry before reading the KV data), but this cost is negligible: the block table is small enough to fit in GPU L2 cache, while the actual KV data dominates the memory access time. In practice, the throughput improvement from higher batch sizes (enabled by reduced waste) far outweighs the minor overhead of the indirection.

\begin{insightbox}
PagedAttention is to GPU memory what virtual memory is to RAM. Just as an OS lets processes see contiguous address spaces backed by scattered physical pages, PagedAttention lets sequences see contiguous KV caches backed by scattered GPU memory blocks. The result is near-zero waste: memory utilization jumps from 20--40\% to over 95\%. This increase in utilization translates directly to higher batch sizes and therefore higher throughput.
\end{insightbox}

\subsection{Quantifying the Improvement}

The impact on serving throughput is dramatic. Consider an A100-80GB serving LLaMA-70B (in a tensor-parallel setup where roughly 10\,GB per GPU is available for the KV cache):

\begin{center}
\begin{tabular}{lrr}
\hline
\textbf{Metric} & \textbf{Without paging} & \textbf{With PagedAttention} \\
\hline
Memory utilization & $\sim$20--40\% & $>$95\% \\
Effective cache capacity & 2--4 GB & 9.5+ GB \\
Max concurrent seqs (4K tokens) & 2--3 & 7--8 \\
Max concurrent seqs (512 tokens) & 6--12 & 59+ \\
\hline
\end{tabular}
\end{center}

For short sequences, the improvement is even more pronounced because the reservation waste under the na\"ive scheme is proportionally larger: a 512-token sequence in a 4,096-token slot wastes 87.5\% of its allocation. PagedAttention reclaims essentially all of that waste.

\subsection{KV Cache Quantization}

\marginnote{\emph{Quantization:} representing values with fewer bits, trading precision for reduced memory and bandwidth.}%
An orthogonal technique for reducing KV cache memory is \emph{quantization}: storing keys and values in lower precision than the model's native format. While the model may compute in FP16 or BF16 (2 bytes per element), the cached keys and values can be quantized to INT8 (1 byte) or even INT4 (0.5 bytes) with minimal impact on output quality.

For LLaMA-70B, moving from FP16 to INT8 halves the per-token cache cost from 320\,KB to 160\,KB. Combined with PagedAttention, this doubles the effective batch size yet again. The trade-off is that quantization introduces small numerical errors in the attention computation---unlike PagedAttention, it is an approximation. However, because keys and values are used only in dot products and weighted sums (not in the more numerically sensitive softmax normalization), the impact on generation quality is typically small. KV cache quantization is supported by most modern inference engines and can be combined with both PagedAttention and prefix sharing.

\section{Prefix Sharing}
\label{sec:kvcache:prefix}

\marginnote{\emph{Prefix sharing:} reusing cached KV blocks across requests that share a common prompt prefix, avoiding redundant computation and storage.}%
PagedAttention eliminates waste \emph{within} a single sequence's allocation. But there is a second, equally large source of waste: \emph{duplication across sequences}. Many production LLM workloads share common prefixes. A system prompt (``You are a helpful assistant\ldots'') may be prepended to every user request. In a coding assistant, the same file context is sent with every completion request. A retrieval-augmented generation (RAG) pipeline may prepend the same set of retrieved documents to multiple queries. Without prefix sharing, each request independently computes and stores KV cache entries for the shared prefix---duplicating both compute and memory.

\subsection{The Opportunity}

Suppose a system prompt of 500 tokens is prepended to every user request. In a batch of 32 concurrent requests, the na\"ive approach stores 32 copies of the same 500-token KV cache---wasting $31 \times 500 \times 320\;\text{KB} \approx 4.8$\,GB for LLaMA-70B. Prefix sharing stores the system prompt's KV cache \emph{once} and lets all 32 sequences reference the same physical blocks:
\begin{equation}
\text{Savings} = (B - 1) \times n_{\text{prefix}} \times c_{\text{token}}
\end{equation}
where $B$ is the batch size, $n_{\text{prefix}}$ is the shared prefix length, and $c_{\text{token}}$ is the per-token cache cost. For our example: $(32 - 1) \times 500 \times 320\;\text{KB} = 4.8$\,GB saved. This is nearly half of the available 10\,GB cache budget, freeing space for additional concurrent requests.

The compute savings are equally significant. Instead of running prefill for the 500-token system prompt 32 times, the system runs it once and reuses the cached keys and values. For a model like LLaMA-70B, prefilling 500 tokens involves roughly $500 \times 2 \times 80 \times 8192^2 \approx 5.4 \times 10^{13}$ FLOPs (considering attention and MLP computations across all layers). Doing this once instead of 32 times saves over $1.6 \times 10^{15}$ FLOPs---multiple seconds of A100 compute time.

\subsection{Mechanism: Block Hashing and Reference Counting}

Prefix sharing requires two capabilities:
\begin{enumerate}
  \item A way to \textbf{identify} whether a given token sequence's KV cache already exists somewhere.
  \item A way to \textbf{share} physical blocks across sequences without copying.
\end{enumerate}

\textbf{Identification via hashing.}\hspace{1em} Token sequences are divided into fixed-size blocks (the same blocks used by PagedAttention), and each block is hashed. In Dynamo, this hashing is performed by the \texttt{dynamo\_tokens} library (Chapter~\ref{ch:tokens}), which computes XXH3 hashes of token blocks:\footnote{XXH3 is a non-cryptographic hash function optimized for speed, producing 64-bit hashes. It runs at over 30\,GB/s on modern CPUs.}
\begin{itemize}
  \item A \emph{local block hash} is computed from the tokens within a single block.
  \item A \emph{sequence hash} chains the local block hash with the previous block's sequence hash, producing a rolling hash that captures the entire prefix up to and including the current block.
\end{itemize}
The sequence hash is critical: two blocks with identical tokens but different preceding contexts will produce different sequence hashes, correctly distinguishing ``\texttt{The cat sat}'' at the start of a prompt from ``\texttt{The cat sat}'' appearing later after different context.

\textbf{Sharing via block pointers.}\hspace{1em} Because PagedAttention already uses block tables with pointers to physical blocks (not contiguous memory), sharing blocks across sequences is straightforward: two sequences' block tables simply point to the same physical blocks for their shared prefix. Reference counting tracks how many sequences use each block. When the last sequence using a shared block finishes, the reference count drops to zero and the block is returned to the free pool.

\begin{figure}[ht]
\centering
\begin{tikzpicture}[
    blocktable/.style={rectangle, draw, minimum width=1.4cm, minimum height=0.5cm, align=center, font=\scriptsize},
    physblock/.style={rectangle, draw, minimum width=1.8cm, minimum height=0.5cm, align=center, font=\scriptsize},
    shared/.style={physblock, fill=green!20},
    unique/.style={physblock, fill=orange!20},
    arrow/.style={-{Stealth[length=4pt]}, thick},
    label/.style={font=\footnotesize\itshape},
    brace/.style={decorate, decoration={brace, amplitude=5pt, mirror}},
]
    % Sequence A block table
    \node[font=\small\bfseries] at (-3.2, 3.5) {Seq.\ A};
    \node[blocktable] (a0) at (-3.2, 2.8) {L0};
    \node[blocktable] (a1) at (-3.2, 2.2) {L1};
    \node[blocktable] (a2) at (-3.2, 1.6) {L2};
    \node[blocktable] (a3) at (-3.2, 1.0) {L3};

    % Sequence B block table
    \node[font=\small\bfseries] at (-3.2, -0.1) {Seq.\ B};
    \node[blocktable] (b0) at (-3.2, -0.8) {L0};
    \node[blocktable] (b1) at (-3.2, -1.4) {L1};
    \node[blocktable] (b2) at (-3.2, -2.0) {L2};
    \node[blocktable] (b3) at (-3.2, -2.6) {L3};

    % Shared physical blocks (system prompt)
    \node[font=\small\bfseries] at (1.5, 3.5) {Physical Blocks};
    \node[shared] (s0) at (1.5, 2.8) {System prompt 0--15};
    \node[shared] (s1) at (1.5, 2.2) {System prompt 16--31};

    % Unique physical blocks for A
    \node[unique] (ua0) at (4.5, 1.6) {A: user tok 32--47};
    \node[unique] (ua1) at (4.5, 1.0) {A: user tok 48--55};

    % Unique physical blocks for B
    \node[unique] (ub0) at (4.5, -2.0) {B: user tok 32--44};
    \node[unique] (ub1) at (4.5, -2.6) {B: user tok 45--51};

    % Arrows from A to shared
    \draw[arrow, green!50!black] (a0.east) -- (s0.west);
    \draw[arrow, green!50!black] (a1.east) -- (s1.west);
    % Arrows from A to unique
    \draw[arrow, orange!70!black] (a2.east) -- (ua0.west);
    \draw[arrow, orange!70!black] (a3.east) -- (ua1.west);

    % Arrows from B to shared (same blocks!)
    \draw[arrow, green!50!black] (b0.east) -- (s0.south west);
    \draw[arrow, green!50!black] (b1.east) -- (s1.south west);
    % Arrows from B to unique
    \draw[arrow, orange!70!black] (b2.east) -- (ub0.west);
    \draw[arrow, orange!70!black] (b3.east) -- (ub1.west);

    % Braces and labels
    \draw[brace] (0.3, 2.0) -- (0.3, 3.1) node[midway, left=6pt, label] {shared (ref count = 2)};
    \draw[brace] (6.6, 0.7) -- (6.6, 1.9) node[midway, right=6pt, label] {A only};
    \draw[brace] (6.6, -2.9) -- (6.6, -1.7) node[midway, right=6pt, label] {B only};
\end{tikzpicture}
\caption{Prefix sharing between two sequences. Sequences A and B share a 32-token system prompt (green blocks). Both block tables point to the \emph{same} physical blocks for the shared prefix; each sequence has its own blocks (orange) only for its unique user tokens. The shared blocks carry a reference count of 2 and are freed only when both sequences complete.}
\label{fig:prefix-sharing}
\end{figure}

\begin{notebox}
\marginnote{\emph{Copy-on-write (CoW):} a strategy where shared data is only duplicated when one user attempts to modify it. Unnecessary for KV caches because cached entries are never modified.}%
Prefix sharing interacts with PagedAttention in an elegant way. PagedAttention was designed to decouple logical and physical memory layout, and prefix sharing is a natural extension: two logical sequences share the same physical blocks for their common prefix, diverging only where their tokens diverge. Copy-on-write semantics could handle the case where a shared block needs modification, but in practice this never occurs for KV caches---once computed, past keys and values are \emph{immutable} (the model weights do not change between tokens, and self-attention is causal, so past positions cannot be influenced by future tokens).
\end{notebox}

\subsection{Radix Trees for Prefix Tracking}

The natural data structure for tracking shared prefixes is a \emph{radix tree} (also known as a Patricia trie or compressed prefix tree). In a radix tree, each path from the root to a node represents a sequence of block hashes---equivalently, a token prefix. Nodes that share a common prefix share the same path in the tree, branching only where the token sequences diverge.

\marginnote{\emph{Radix tree:} a compressed trie where each edge represents a sequence of elements (here, token block hashes). Shared prefixes share edges, branching only at points of divergence.}%
Dynamo implements a radix tree in its \texttt{kv-router} library that maps block hash sequences to the set of workers that have those blocks cached. Each node in the tree (a \texttt{RadixBlock} in the code) stores a set of workers that have cached the corresponding prefix, along with child pointers keyed by local block hash. When a new request arrives, the router walks the radix tree using the request's block hashes, finding the longest matching prefix. The worker with the most cached blocks for that prefix is selected to serve the request, maximizing cache reuse and minimizing redundant computation.

The hashing scheme in Dynamo is carefully designed to be both fast and collision-resistant. The \texttt{dynamo\_tokens} library uses a 128-bit \texttt{PositionalSequenceHash} that encodes three pieces of information: the 64-bit rolling sequence hash (which captures the full prefix history), the block's position in the sequence, and the local block hash. The position encoding uses a variable-width scheme---8, 16, 24, or 31 bits depending on the sequence length---to maximize the hash bits available for collision resistance at typical context lengths.

This is the direct connection between the KV cache concepts in this chapter and the routing system in Chapter~\ref{ch:kvrouter}: the radix tree exists \emph{because} KV cache blocks are immutable, hashable, and shareable. The entire KV-aware routing strategy is built on the observation that prefix sharing turns cache placement into a routing decision. Where a traditional load balancer distributes requests based on CPU utilization or connection count, Dynamo's KV-aware router distributes requests based on \emph{which GPU already has the right KV cache blocks in memory}---a fundamentally different optimization target.

\begin{insightbox}
Prefix sharing transforms KV cache management from a per-request problem into a \emph{system-wide} optimization. Instead of each request independently managing its own cache, the system maintains a shared pool of cached prefixes and routes requests to maximize reuse. This is why Dynamo's router needs to know which GPU has which cached blocks---it is implementing prefix sharing at the \emph{cluster} level, not just within a single GPU's memory.
\end{insightbox}

\begin{whynotbox}{Why not cache intermediate activations too?}
If caching keys and values avoids redundant work, why stop there? The transformer also computes intermediate activations at every layer---the residual stream, attention outputs, FFN outputs. Could we cache those too?

\textbf{Activations are not reusable across sequences.} A key vector $\mathbf{k}_j = \mathbf{W}_K \mathbf{x}_j$ depends only on position $j$'s input and the model weights. It is read (via dot products) by every future token's query, making it valuable to cache. An intermediate activation like the FFN output at position $j$ in layer $\ell$, by contrast, is consumed exactly once---by the residual connection feeding layer $\ell + 1$---and then never needed again during inference. Caching a value that is used once and discarded has no benefit.

\textbf{Keys and values have a unique access pattern.} During decode, the attention mechanism at each step reads \emph{all previous} keys and values (the full $\mathbf{K}_{1:t}$ and $\mathbf{V}_{1:t}$). This is the only part of the forward pass that accesses prior positions' data. All other computations (layer norms, FFN, projections) operate only on the \emph{current} token's representation. The KV cache exists precisely because attention is the one operation that requires cross-position data from prior steps.

\textbf{Memory cost would be prohibitive.} Caching all intermediate activations for a sequence of $t$ tokens across $L$ layers would require $O(t \times L \times d)$ memory---roughly $L \times$ the KV cache. For LLaMA-70B with 80 layers, this would multiply cache memory by a factor of 40, with zero benefit for inference speed.
\end{whynotbox}

\subsection{Limitations and Trade-offs}

Prefix sharing is not free. The hash computation adds overhead (though XXH3 is extremely fast). The radix tree must be maintained as workers evict blocks under memory pressure, requiring a protocol for cache invalidation events. And prefix sharing only helps when requests actually share prefixes---in a workload with entirely unique prompts, there is nothing to share.

In practice, most production workloads have substantial prefix overlap. System prompts, few-shot examples, tool-use preambles, and RAG-injected context all create shared prefixes ranging from hundreds to thousands of tokens. The savings compound with batch size: the more concurrent requests there are, the more likely it is that several share a prefix, and the more memory and compute are saved.

\marginnote{\emph{Eviction:} removing cached blocks from GPU memory to make room for new sequences. The eviction policy determines which blocks to sacrifice.}%
There is also a subtlety around \emph{eviction}. When a GPU runs low on free blocks, it must evict some cached blocks to make room for new sequences. A na\"ive LRU (least recently used) eviction policy might evict blocks from a popular shared prefix---destroying the cache for many sequences simultaneously. Smarter eviction policies account for reference counts and sharing patterns: a block shared by 10 active sequences is far more valuable than a block used by a single completed sequence, even if the latter was accessed more recently. Dynamo's radix tree tracks per-block timestamps (the \texttt{recent\_uses} field) and worker sets, providing the data needed for informed eviction decisions.

Another challenge is \emph{cache coherence} in a distributed setting. When a worker evicts blocks under memory pressure, the central radix tree must be updated to reflect that those blocks are no longer available on that worker. Stale entries---where the tree believes a worker has a block that has actually been evicted---lead to suboptimal routing decisions: the router sends a request to a worker expecting a cache hit, only to discover the blocks are gone, forcing a full recomputation. Dynamo addresses this with an event-based protocol where workers emit cache events (stores and removes) that are consumed by the KV indexer to keep the radix tree in sync (Chapter~\ref{ch:kvrouter}).

\section{Summary}

\begin{itemize}
  \item The \textbf{KV cache} stores precomputed key and value vectors to avoid redundant computation during autoregressive generation, trading memory for a factor-of-$T$ reduction in projection compute.
  \item For each layer and KV head, the cache stores a key matrix and value matrix that grow by one row per decode step. With GQA, the number of KV heads ($H_{\text{KV}}$) is much smaller than the number of query heads ($H$), directly reducing cache size.
  \item Cache memory per token is $2 \times L \times H_{\text{KV}} \times d_{\text{head}} \times b$ bytes, and total cache memory grows as $O(B \times s \times L \times H_{\text{KV}} \times d_{\text{head}})$---\textbf{dominating GPU memory} for large batches with long sequences.
  \item \textbf{PagedAttention} applies OS-style paging to eliminate memory fragmentation, increasing utilization from 20--40\% to over 95\%. Memory is divided into fixed-size blocks allocated on demand, tracked via per-sequence block tables.
  \item \textbf{Prefix sharing} allows requests with common prefixes to reuse cached KV blocks, saving both compute and memory. Identification uses block hashing; sharing uses reference-counted block pointers.
  \item A \textbf{radix tree} tracks which prefixes are cached on which workers, enabling KV-aware routing at the cluster level---the subject of Chapter~\ref{ch:kvrouter}.
  \item \textbf{KV cache quantization} (INT8, INT4) is an orthogonal technique that halves or quarters the per-token cache cost at the expense of small numerical approximation.
  \item The \textbf{prefill-decode asymmetry}---prefill is compute-bound, decode is memory-bound---motivates disaggregated serving, which in turn requires efficient KV cache transfer between GPUs.
  \item PagedAttention and prefix sharing are not approximations---they produce \textbf{bit-identical} results to standard attention while dramatically improving resource efficiency.
\end{itemize}

The key equations to remember:
\begin{align*}
\text{Per-token cache} &= 2 \times L \times H_{\text{KV}} \times d_{\text{head}} \times b \\
\text{Total cache} &= B \times s \times \text{(per-token cache)} \\
\text{Na\"ive projection work} &= O(nT + T^2) \\
\text{Cached projection work} &= O(n + T)
\end{align*}

\medskip
\begin{center}\rule{0.5\linewidth}{0.4pt}\end{center}
\medskip

With the KV cache understood, we can now ask: what happens when a single GPU's memory is not enough to hold the model, the cache, and the batch? The next chapter answers this question by introducing distributed inference---tensor parallelism, pipeline parallelism, and the disaggregation of prefill and decode.

\medskip
\noindent\textbf{Check Your Understanding}

\smallskip
\noindent\emph{These questions start with recall and progress to conceptual reasoning. All are answerable from this chapter alone.}

\begin{enumerate}
  \item Why is the KV cache necessary? What would happen during autoregressive generation without it, and what is the asymptotic difference in total projection work?

  \item For a model with 32 layers, 8 KV heads, $d_k = 128$, and FP16 precision, how many bytes of KV cache does a single token require? Show your calculation.

  \item Explain the three forms of memory waste that PagedAttention eliminates. Which form is typically the most severe, and why?

  \item Why is prefix sharing particularly valuable for production LLM workloads? Give three concrete scenarios where it applies.

  \item PagedAttention and prefix sharing both improve memory efficiency, but they address different problems. What specific inefficiency does each one target?

  \item A system serves LLaMA-70B on 4 A100-80GB GPUs with tensor parallelism. After loading model weights, each GPU has approximately 10\,GB free for the KV cache. What is the maximum number of concurrent 4,096-token sequences the system can support? What if the average sequence is only 512 tokens and PagedAttention is used?

  \item Consider a chatbot deployment where every request starts with a 1,000-token system prompt. The system serves a batch of 64 concurrent requests on a GPU with 10\,GB of cache budget. How much memory does prefix sharing save compared to the na\"ive approach? Express your answer both in absolute terms and as a percentage of the total cache budget.

  \item The chapter states that KV cache entries are immutable---once computed, a position's key and value vectors never change. Why is this property essential for prefix sharing to be correct? What would go wrong if keys or values could be modified after computation?
\end{enumerate}

\medskip
\noindent\textbf{Answers}

\begin{enumerate}
  \item Without the KV cache, every decode step must recompute key and value projections for all previous tokens. Generating $T$ tokens from an $n$-token prompt requires $O(nT + T^2)$ total projection sets across all steps (Equation~\ref{eq:naive-projections}). With the cache, each step computes only the new token's projections, reducing total projection work to $O(n + T)$ (Equation~\ref{eq:cached-projections})---a factor-of-$T$ improvement. The attention computation itself ($O(T^2 d)$ total) is unchanged, since attending to all previous positions is inherent to the algorithm.

  \item Using Equation~\ref{eq:per-token}: $2 \times 32 \times 8 \times 128 \times 2 = 131{,}072$ bytes $= 128$\,KB per token. The factor of 2 in front accounts for storing both keys and values; the trailing factor of 2 accounts for FP16's 2 bytes per element.

  \item (1)~\textbf{Reservation waste}: pre-allocating for maximum sequence length when most sequences are shorter; (2)~\textbf{Internal fragmentation}: allocated buffers exceeding the actual final sequence length; (3)~\textbf{External fragmentation}: free memory scattered into non-contiguous chunks too small to serve new requests. Reservation waste is typically the most severe because it scales with the ratio of maximum to average sequence length. A system with a 4,096-token maximum serving sequences that average 200 tokens wastes over 95\% of allocated cache memory.

  \item (1)~System prompts (``You are a helpful assistant\ldots'') prepended to every request; (2)~few-shot examples included in every prompt for a specific task; (3)~RAG pipelines that inject the same retrieved documents for multiple user queries about the same topic. Each scenario creates a prefix of hundreds to thousands of tokens shared across many concurrent requests.

  \item PagedAttention targets \emph{fragmentation}---the waste from pre-allocating contiguous memory at maximum sequence length and from scattered free blocks. Prefix sharing targets \emph{duplication}---the waste from storing identical KV cache entries multiple times when requests share a common prefix. They are complementary: PagedAttention reduces waste within a single sequence's allocation, while prefix sharing reduces waste across sequences.

  \item Per token: $2 \times 80 \times 8 \times 128 \times 2 = 327{,}680$ bytes $\approx 320$\,KB. With 10\,GB free: $10 \times 1024^3 / (320 \times 1024) \approx 32{,}768$ token-slots. At 4,096 tokens per sequence: $32{,}768 / 4{,}096 = 8$ concurrent sequences per GPU, or 32 total across 4 GPUs. With 512-token sequences and PagedAttention (so memory is allocated only for actual tokens, not reserved for 4,096): $32{,}768 / 512 = 64$ concurrent sequences per GPU, or 256 total. PagedAttention's on-demand allocation makes this possible; without it, each 512-token sequence would still reserve 4,096 tokens of space, limiting throughput to the same 8 per GPU.

  \item Without prefix sharing: 64 sequences $\times$ (1,000 prefix tokens + user tokens). Focusing just on the prefix: $64 \times 1{,}000 \times 320$\,KB $= 20{,}480$\,MB $\approx 20$\,GB for the prefix alone (which already exceeds the 10\,GB budget---the system would need to reduce batch size). With prefix sharing: $1 \times 1{,}000 \times 320$\,KB $= 320$\,MB for one copy. Savings: $63 \times 1{,}000 \times 320$\,KB $= 19.7$\,GB. As a percentage of the 10\,GB budget, the na\"ive approach requires 200\% of the budget for prefixes alone; prefix sharing reduces this to 3.1\%. The savings are 197\% of the cache budget---meaning prefix sharing is the difference between the workload being impossible and feasible.

  \item Immutability is essential because prefix sharing allows multiple sequences to read from the same physical memory blocks. If one sequence could modify a shared key or value vector, it would corrupt the cache for all other sequences sharing that block---a classic shared-mutable-state problem. Immutability guarantees that sharing is safe without locks or copy-on-write: once a block is computed, its contents are fixed, so any number of sequences can read from it concurrently with no risk of data races or corruption. If keys or values could change (for example, in a hypothetical architecture where attention weights are updated based on later context), prefix sharing would require either copying blocks before modification (negating the memory savings) or complex locking protocols (adding latency and deadlock risk).
\end{enumerate}
